% Replacement for Chapter 19. This file includes its own \chapter command.
% Integrate it in place of the old chapter heading and 19.1--19.4 inputs.
\chapter{States and Correlations}
\label{ch:correlations}

Chapter~18 showed how a hyperbolic change of current variables simplifies
an interacting quadratic Hamiltonian.  A diagonal expression suggests a
vacuum and, with it, correlation functions.  There is a difficulty, however:
the finite fermionic system and the familiar bosonic oscillator have
different algebraic properties.  Their currents agree on appropriately
protected low-energy vectors, but they are not globally the same operators.
Which vacuum, then, is being used when we calculate a correlation?

We will give two precise answers.  First, an \emph{uncompressed oscillator
algebra} has a positive state in which the dressed modes annihilate the
vacuum.  Its correlations follow by finite algebraic calculations.  Second,
the original finite CAR system has genuine vector states, and a protected
fermionic sea permits some of the same current calculations.  Their answers
agree in the free case; an interacting-state comparison requires additional
work.  The distinction will remain visible in the notation throughout.

\section{Two kinds of quantum expectation}
\label{sec:ch19-scope}

\subsection*{The finite system}

Let $\mathcal H_{\mathrm F}$ be the finite-dimensional fermionic Fock space
constructed earlier.  A normalized vector $\Psi\in\mathcal H_{\mathrm F}$
and a positive trace-one matrix $\Gamma$ give, respectively,
\begin{equation}
 \langle X\rangle^{\mathrm F}_{\Psi}
 :=\langle\Psi,X\Psi\rangle,
 \qquad
 \langle X\rangle^{\mathrm F}_{\Gamma}
 :=\operatorname{Tr}(\Gamma X).
 \label{eq:ch19-vector-expectation}
\end{equation}
The inner product is linear in the second argument.  Both expressions are
normalized positive linear functionals: in particular,
$\langle X^*X\rangle^{\mathrm F}_{\Psi}=\|X\Psi\|^2\geq0$.
For a specified self-adjoint finite Hamiltonian, a ground-state vector
always exists.  It need not be annihilated by operators called
``dressed annihilators.''  The latter is a separate condition.

The value of an ordered product depends on its order:
$\langle XY\rangle$ and $\langle YX\rangle$ may differ, even in a
ground state.  Neither is generally the product of the separate
expectations.  These elementary observations will determine the two-point
functions below.

\subsection*{Why the finite currents are not global oscillators}

A one-line argument marks the boundary of the finite description.

\begin{proposition}[Finite-dimensional trace obstruction]
\label{prop:ch19-trace-obstruction}
If $A,C$ are endomorphisms of a nonzero finite-dimensional complex
vector space and $[A,C]=\lambda I$, then $\lambda=0$.
\end{proposition}
\begin{proof}
Cyclicity of the matrix trace gives
$0=\operatorname{Tr}(AC-CA)=\lambda\dim\mathcal H$.
\end{proof}

In particular, the canonical commutator $[A,C]=mI$, $m>0$, cannot hold
throughout finite Fock space.  This does \emph{not} contradict the finite
Schwinger-term identities proved in earlier chapters: those identities
act on selected vectors under band-margin hypotheses, and they do not
assert a global scalar commutator.

There is a second way to see the issue.  For an orthogonal projection
$\Pi$ and $Q=I-\Pi$, compression gives
\begin{equation}
 [\Pi A\Pi,\Pi C\Pi]
 =\Pi[A,C]\Pi-\Pi AQ C\Pi+\Pi CQ A\Pi,
 \label{eq:ch19-compression-remainder}
\end{equation}
where the equality is understood on $\Pi\mathcal H$.  The last two
terms measure paths that leave the projected space between the two
operator actions.  They need not vanish simply because the initial
and final vectors lie below an energy cutoff.

\subsection*{Finite-band boundary currents and the exact boundary correction}

To make the microscopic mechanics of the finite CAR system transparent,
consider a finite momentum band $\mathcal B_B = \{-B, -B+1, \ldots, B\}$
of half-width $B \geq 1$.  Fermionic creation and annihilation operators
$c_k^\dagger, c_k$ satisfy the canonical anticommutation relations
$\{c_k, c_p\} = 0$ and $\{c_k, c_p^\dagger\} = \delta_{kp} I$.
For indices $i, j \in \mathcal B_B$, define the elementary bilinear operators
\begin{equation}
 E_{ij} := c_i^\dagger c_j.
 \label{eq:ch19-bilinear-def}
\end{equation}

\begin{proposition}[Bilinear Fermionic Commutator]
\label{prop:ch19-bilinear-commutator}
For any four orbitals $i, j, k, l \in \mathcal B_B$, the bilinears satisfy
\begin{equation}
 [E_{ij}, E_{kl}] = \delta_{jk} E_{il} - \delta_{li} E_{kj}.
 \label{eq:ch19-bilinear-comm}
\end{equation}
\end{proposition}
\begin{proof}
Expand the commutator directly from $E_{ij} = c_i^\dagger c_j$:
\[
 [E_{ij}, E_{kl}] = c_i^\dagger c_j c_k^\dagger c_l - c_k^\dagger c_l c_i^\dagger c_j.
\]
Using the CAR identities $c_j c_k^\dagger = \delta_{jk}I - c_k^\dagger c_j$ and
$c_l c_i^\dagger = \delta_{li}I - c_i^\dagger c_l$:
\begin{align*}
 [E_{ij}, E_{kl}]
   &= c_i^\dagger (\delta_{jk}I - c_k^\dagger c_j) c_l - c_k^\dagger (\delta_{li}I - c_i^\dagger c_l) c_j \\
   &= \delta_{jk} c_i^\dagger c_l - \delta_{li} c_k^\dagger c_j
      - c_i^\dagger c_k^\dagger c_j c_l + c_k^\dagger c_i^\dagger c_l c_j.
\end{align*}
By anticommutativity of fermionic operators, $c_i^\dagger c_k^\dagger = -c_k^\dagger c_i^\dagger$
and $c_j c_l = -c_l c_j$.  Therefore
\[
 - c_i^\dagger c_k^\dagger c_j c_l + c_k^\dagger c_i^\dagger c_l c_j
  = c_k^\dagger c_i^\dagger (-c_l c_j) + c_k^\dagger c_i^\dagger c_l c_j = 0.
\]
The quartic fermion terms cancel identically, leaving $\delta_{jk} E_{il} - \delta_{li} E_{kj}$.
\end{proof}

For a positive integer displacement $n \in \{1, \ldots, B\}$, the finite-band current
operator is the non-wrapping bilinear sum
\begin{equation}
 J_n^{(B)} := \sum_{k=-B}^{B-n} c_k^\dagger c_{k+n} = \sum_{k=-B}^{B-n} E_{k, k+n},
 \qquad
 J_{-n}^{(B)} := (J_n^{(B)})^\dagger = \sum_{l=-B}^{B-n} c_{l+n}^\dagger c_l = \sum_{l=-B}^{B-n} E_{l+n, l}.
 \label{eq:ch19-finite-current-def}
\end{equation}
Because the band is finite, this sum terminates strictly at $B-n$ rather than running
to infinity.  The boundary of the band produces an exact correction to the canonical
current commutator:

\begin{theorem}[Exact Boundary-Corrected Finite Current Commutator]
\label{thm:ch19-current-boundary-commutator}
For any $B \ge 1$ and $1 \le n \le B$, the finite-band currents obey the exact operator identity
\begin{equation}
 \boxed{[J_n^{(B)}, J_{-n}^{(B)}]
  = \sum_{k=-B}^{-B+n-1} \hat{n}_k - \sum_{k=B-n+1}^B \hat{n}_k,}
 \label{eq:ch19-current-exact-boundary}
\end{equation}
where $\hat{n}_k = c_k^\dagger c_k = E_{kk}$ is the number operator on orbital $k$.
\end{theorem}
\begin{proof}
Expand the commutator of the two finite sums using Proposition~\ref{prop:ch19-bilinear-commutator}:
\begin{align*}
 [J_n^{(B)}, J_{-n}^{(B)}]
  &= \sum_{k=-B}^{B-n} \sum_{l=-B}^{B-n} [E_{k, k+n}, E_{l+n, l}] \\
  &= \sum_{k=-B}^{B-n} \sum_{l=-B}^{B-n} \left(\delta_{k+n, l+n} E_{k, l} - \delta_{k, l} E_{l+n, k+n}\right).
\end{align*}
In the first term, $\delta_{k+n, l+n} = \delta_{kl}$, yielding $E_{kk} = \hat{n}_k$.
In the second term, $\delta_{kl}$ sets $l = k$, yielding $E_{k+n, k+n} = \hat{n}_{k+n}$.
Therefore:
\[
 [J_n^{(B)}, J_{-n}^{(B)}]
  = \sum_{k=-B}^{B-n} (\hat{n}_k - \hat{n}_{k+n}).
\]
Reindex the second sum by setting $j = k+n$:
\[
 \sum_{k=-B}^{B-n} \hat{n}_{k+n} = \sum_{j=-B+n}^B \hat{n}_j.
\]
For $1 \le n \le B$, the two index sets $\{-B, \ldots, B-n\}$ and $\{-B+n, \ldots, B\}$
overlap on the interior range $\{-B+n, \ldots, B-n\}$.
Subtracting the two sums causes the interior terms to cancel completely:
\[
 \sum_{k=-B}^{B-n} \hat{n}_k - \sum_{j=-B+n}^B \hat{n}_j
  = \sum_{k=-B}^{-B+n-1} \hat{n}_k - \sum_{k=B-n+1}^B \hat{n}_k.
\]
This completes the proof.
\end{proof}

\begin{corollary}[Action on Margin-Protected States]
\label{cor:ch19-margin-protected-action}
Let $|S\rangle$ be an occupation basis configuration.  If every mode in the bottom margin
$\{-B, \ldots, -B+n-1\}$ is occupied ($\hat{n}_k|S\rangle = |S\rangle$) and every mode in the
top margin $\{B-n+1, \ldots, B\}$ is empty ($\hat{n}_k|S\rangle = 0$), then
\begin{equation}
 [J_n^{(B)}, J_{-n}^{(B)}]|S\rangle = n |S\rangle.
 \label{eq:ch19-protected-eigenvalue}
\end{equation}
Conversely, if an excitation empties an orbital in the bottom margin or occupies an orbital
in the top margin, the eigenvalue deviates strictly from $n$.
\end{corollary}

This makes the architectural boundary vivid: the finite-band current commutator acts as the
scalar $n\cdot I$ \emph{only on states within a certified margin budget}.  If one defines a
projected operator $\Pi J_n \Pi$, compression destroys the algebra by
\eqref{eq:ch19-compression-remainder}; uncompressed products must be evaluated on protected states.

We therefore distinguish the finite expectation
$\langle\,\cdot\,\rangle^{\mathrm F}$ from the abstract bosonic
functional $\omega_P$ constructed next.  Here $P=(c,s)$ denotes
Bogoliubov parameters; $\Pi$ is reserved for a projection.

\section{Building the dressed oscillator state}
\label{sec:ch19-oscillator-state}

\subsection*{An algebra with no highest occupation}

Fix a finite set of positive modes $m=1,\ldots,M$ and two branches
$\nu\in\{R,L\}$; $M=0$ is allowed, with empty sums understood.  For each
label introduce bare creation and annihilation generators $C_{\nu m}$
and $A_{\nu m}=C_{\nu m}^{*}$.  They generate the unital complex
$*$-algebra $\mathfrak A_M$ subject to
\begin{equation}
 [A_{\nu m},C_{\eta n}]
   =m\delta_{\nu\eta}\delta_{mn}I,
 \qquad [A_{\nu m},A_{\eta n}]
   =[C_{\nu m},C_{\eta n}]=0.
 \label{eq:ch19-abstract-ccr}
\end{equation}
Here $*$ reverses the order of products as well as taking complex
conjugates of coefficients.  Algebraically, $\mathfrak A_M$ can be
constructed from finite noncommuting words by identifying expressions
that differ by the relations in \eqref{eq:ch19-abstract-ccr}.  The
identification is compatible with multiplication on both sides: an
identity remains valid when inserted into a longer word.  No
infinite series or limiting procedure is involved.

It is useful to see an actual space on which these relations hold.
Consider commuting polynomial variables $X_{\nu m}$ and the carrier
\[
 \mathcal P_M=\mathbb C[X_{\nu m}:\nu\in\{R,L\},\ 1\leq m\leq M].
\]
The \emph{variables} commute; their multiplication and differentiation
\emph{operators} do not.  Define
\begin{equation}
 \pi(C_{\nu m})f=X_{\nu m}f,
 \qquad \pi(A_{\nu m})f=m\frac{\partial f}{\partial X_{\nu m}},
 \label{eq:ch19-polynomial-representation}
\end{equation}
and make the monomials orthogonal with
\begin{equation}
 \langle X^{\alpha},X^{\beta}\rangle_{\mathcal P}
  =\delta_{\alpha\beta}
   \prod_{\nu,m}\alpha_{\nu m}!\,m^{\alpha_{\nu m}}.
 \label{eq:ch19-polynomial-gram}
\end{equation}
The positive weights ensure a genuine inner product on polynomials,
without taking a Hilbert-space completion.  Multiplication by
$X_{\nu m}$ and $m\partial/\partial X_{\nu m}$ are adjoints on this
common polynomial domain, and the ordinary differentiation rule proves
\eqref{eq:ch19-abstract-ccr}.  Thus these operators represent the
relation algebra consistently and show it is nonzero.

The smallest example already captures the mechanism.  With one mode of
weight $m$, let $C=X$ and $A=m\,d/dX$.  Then
\begin{equation}
 AC\,X^n=m(n+1)X^n,
 \qquad CA\,X^n=mnX^n,
 \label{eq:ch19-one-mode-ccr}
\end{equation}
so $[A,C]=mI$ for every $n\geq0$.  Each polynomial has finite
degree, but there is no highest allowed degree.  If one instead
retains only $1,X,\ldots,X^N$ and compresses the operators, their
commutator acts as $mI$ below the top level and as $-mN$ on $X^N$.
The defect at that one level restores the zero trace required in
finite dimension.  This is why an energy-truncated oscillator is not
an exact uncompressed oscillator.

\subsection*{A positive vacuum without an infinite series}

Write $\mathbf1$ for the constant polynomial, which has norm one.
It is annihilated by every $A_{\nu m}$.  The vacuum expectation is
\begin{definition}[Polynomial vacuum]
\label{def:ch19-polynomial-vacuum}
\begin{equation}
 \omega_0(a):=\langle\mathbf1,\pi(a)\mathbf1\rangle_{\mathcal P},
 \qquad a\in\mathfrak A_M.
 \label{eq:ch19-omega-zero}
\end{equation}
\end{definition}

\begin{proposition}[Positivity and the vacuum rule]
\label{prop:ch19-polynomial-state}
The functional $\omega_0$ is normalized and positive.  For every
$a\in\mathfrak A_M$ and retained mode,
\begin{equation}
 \omega_0(a^*a)=\|\pi(a)\mathbf1\|_{\mathcal P}^2\geq0,
 \qquad
 \omega_0(aA_{\nu m})=\omega_0(C_{\nu m}a)=0.
 \label{eq:ch19-omega0-positive}
\end{equation}
\end{proposition}
\begin{proof}
Normalization follows from $\|\mathbf1\|=1$; positivity follows
from the algebraic adjoint relation.  A rightmost annihilator acts
first and kills $\mathbf1$.  The leftmost-creator identity is its
adjoint counterpart.
\end{proof}

These two rules give the fundamental ordered contractions:
\begin{equation}
 \omega_0(A_{\nu m}C_{\eta n})
     =m\delta_{\nu\eta}\delta_{mn},
 \qquad \omega_0(C_{\nu m}A_{\eta n})=0.
 \label{eq:ch19-free-two-orders}
\end{equation}
For instance, $AC\mathbf1=m\mathbf1$, whereas
$CA\mathbf1=0$.  The nonzero expectation in the first order is a
vacuum fluctuation fixed by the CCR, not a thermodynamic
approximation.  Expectations of longer finite words reduce by
commuting annihilators to the right; this pairing rule is the
algebraic meaning of a Gaussian, or quasi-free, oscillator vacuum.

\subsection*{A vacuum for the dressed modes}

Take real $c,s$ with $c^2-s^2=1$, as in Chapter~18, and let
$\bar R=L$, $\bar L=R$.  The Bogoliubov map on bare generators is
\begin{equation}
 \begin{aligned}
 \beta(C_{\nu m})&=cC_{\nu m}+sA_{\bar\nu m},\\
 \beta(A_{\nu m})&=cA_{\nu m}+sC_{\bar\nu m}.
 \end{aligned}
 \label{eq:ch19-bogoliubov-creator}
\end{equation}
The transformed generators satisfy the same CCR: in particular,
$[\beta(A_{Rm}),\beta(C_{Rn})]
=m\delta_{mn}(c^2-s^2)I$.  The cross-branch terms cancel in the
remaining commutators.  Reality of $c,s$ preserves adjoints, and
replacing $s$ by $-s$ gives the inverse.  Consequently $\beta$ is a
$*$-automorphism of $\mathfrak A_M$.

Denote its images by
$\widetilde C_{\nu m}=\beta(C_{\nu m})$ and
$\widetilde A_{\nu m}=\beta(A_{\nu m})$.  We want a state in which
the \emph{dressed} annihilators, rather than the bare annihilators,
kill the vacuum.  This fixes the direction of the pullback:
\begin{definition}[Dressed abstract state]
\label{def:ch19-dressed-functional}
\begin{equation}
 \boxed{\omega_P(a)=\omega_0\bigl(\beta^{-1}(a)\bigr),
 \qquad P=(c,s).}
 \label{eq:ch19-dressed-state}
\end{equation}
\end{definition}

\begin{proposition}[Positive dressed vacuum]
\label{prop:ch19-dressed-positive}
The functional $\omega_P$ is normalized and positive, and
\begin{equation}
 \omega_P(a\widetilde A_{\nu m})=0,
 \qquad \omega_P(\widetilde C_{\nu m}a)=0
 \qquad (a\in\mathfrak A_M).
 \label{eq:ch19-dressed-annihilation}
\end{equation}
\end{proposition}
\begin{proof}
A $*$-automorphism maps $a^*a$ to the corresponding positive
square, so positivity follows from that of $\omega_0$.  Moreover,
$\beta^{-1}(\widetilde A_{\nu m})=A_{\nu m}$ and
$\beta^{-1}(\widetilde C_{\nu m})=C_{\nu m}$; the preceding vacuum
rules give both vanishing statements.
\end{proof}

This is a \emph{constructed state on an uncompressed polynomial
$*$-algebra}, not a density matrix on finite fermionic Fock space.
In particular, we have neither assumed a global finite CCR nor
constructed an interacting finite-CAR ground ket.

\section{Mode correlations: why order matters}
\label{sec:ch19-mode-correlators}

An observable expressed in the original density modes uses the
\emph{bare} generators.  Although $\omega_P$ is a vacuum for the
dressed modes, it need not annihilate their bare counterparts.
To find its measurable consequences, invert the transformation:
\[
 \beta^{-1}(C_{\nu m})=cC_{\nu m}-sA_{\bar\nu m},
 \qquad
 \beta^{-1}(A_{\nu m})=cA_{\nu m}-sC_{\bar\nu m}.
\]
The following finite table is the essential calculation of the chapter.

\begin{theorem}[Ordered contractions in the dressed state]
\label{thm:ch19-mode-same-branch}
\label{thm:ch19-mode-cross-branch}
For $1\leq m,n\leq M$ and $\nu\in\{R,L\}$,
\begin{align}
 \omega_P(C_{\nu m}A_{\nu n})
   &=m\delta_{mn}s^2,
   \label{eq:ch19-mode-CA}\\
 \omega_P(A_{\nu m}C_{\nu n})
   &=m\delta_{mn}c^2,
   \label{eq:ch19-mode-AC}\\
 \omega_P(C_{\nu m}C_{\nu n})
   &=\omega_P(A_{\nu m}A_{\nu n})=0.
   \label{eq:ch19-mode-like}
\end{align}
For opposite branches,
\begin{align}
 \omega_P(C_{Rm}C_{Ln})
   &=-m\delta_{mn}cs,
   \label{eq:ch19-cross-CC}\\
 \omega_P(A_{Rm}A_{Ln})
   &=-m\delta_{mn}cs,
   \label{eq:ch19-cross-AA}\\
 \omega_P(C_{Rm}A_{Ln})
   &=\omega_P(A_{Rm}C_{Ln})=0.
   \label{eq:ch19-cross-mixed}
\end{align}
The cross-branch relations also hold with $R$ and $L$ exchanged.
\end{theorem}
\begin{proof}
Consider the two terms that reveal the mechanism.  In the same
branch, the inverse pullback gives
\begin{align*}
 \omega_P(C_{Rm}A_{Rn})
 &=\omega_0\!\left((cC_{Rm}-sA_{Lm})
                    (cA_{Rn}-sC_{Ln})\right)
 =s^2m\delta_{mn}.
\end{align*}
Only $A_{Lm}C_{Ln}$ survives the bare vacuum evaluation.  On the
other hand,
\[
 \omega_P(A_{Rm}C_{Rn})
 =\omega_0\!\left((cA_{Rm}-sC_{Lm})
                     (cC_{Rn}-sA_{Ln})\right)
 =c^2m\delta_{mn},
\]
because now $A_{Rm}C_{Rn}$ survives.  For a cross-branch pair,
\[
 \omega_P(C_{Rm}C_{Ln})
 =\omega_0\!\left((cC_{Rm}-sA_{Lm})
                     (cC_{Ln}-sA_{Rn})\right)
 =-cs\,m\delta_{mn}.
\]
The minus sign comes from the inverse map.  The other formulas
follow by the same two-point vacuum rule, adjoints, and commuting
different branches.
\end{proof}

Two consequences deserve emphasis.  First, the difference of the two
same-branch orders is fixed:
\begin{equation}
 \omega_P(A_{\nu m}C_{\nu n})
 -\omega_P(C_{\nu n}A_{\nu m})
 =m\delta_{mn}.
 \label{eq:ch19-mode-commutator-check}
\end{equation}
Replacing both coefficients by $s^2$ would already give the wrong
free-vacuum limit.  Second, opposite-branch generators commute but
have nonzero correlations in $\omega_P$.  Commutativity is not
statistical independence.

The Luttinger parameter $K > 0$ determines the dressing angle $\theta = \frac{1}{2}\ln K$.
The hyperbolic coefficients are uniquely locked by
\begin{equation}
 c = \cosh\theta = \frac{\sqrt{K} + 1/\sqrt{K}}{2}, \qquad
 s = \sinh\theta = \frac{\sqrt{K} - 1/\sqrt{K}}{2}.
 \label{eq:ch19-hyperbolic-locked}
\end{equation}
They satisfy the fundamental algebraic identities:
\begin{equation}
 c^2 - s^2 = 1, \qquad
 c^2 + s^2 = \frac{K + K^{-1}}{2}, \qquad
 2cs = \frac{K - K^{-1}}{2}, \qquad
 \sinh(2\theta) = \frac{K^2 - 1}{2K}.
 \label{eq:ch19-mode-g-coefficients}
\end{equation}
In the notation of Chapter~18, $g = (c-s)^2 = 1/K$ and $g^{-1} = (c+s)^2 = K$.
We explicitly warn the reader that the alternative convention $\sinh(2\theta) = (1-K^2)/(2K)$
corresponds to $\theta = -\frac{1}{2}\ln K$ and is rejected throughout this monograph to maintain
strict internal consistency.
The symmetrized same-branch contraction is therefore
$\tfrac12m\delta_{mn}(c^2+s^2) = \frac{K+K^{-1}}{4}m\delta_{mn}$.  In contrast, the ordered
products retain the separate $c^2$ and $s^2$ coefficients.

This is the fluctuation content of the dressing.  Relative to the
dressed vacuum, a bare mode has occupation $s^2$ per unit mode weight,
and the two branches acquire an anomalous pair correlation $-cs$.
Thus mixing does not remove vacuum fluctuations; it redistributes
them between branches.  The same-branch symmetric fluctuation grows
from its free value by the factor $c^2+s^2=1+2s^2$, while its
antisymmetric part remains fixed by the CCR.  These are exact
covariances of the abstract state defined here, not a thermal
population or a claim about a finite interacting fermionic ground state.

As a quick positivity check, fix $m$ and put
$Y_m=C_{Rm}+A_{Lm}$.  The theorem gives
\begin{equation}
 \omega_P(Y_m^*Y_m)
   =m(c^2+s^2-2cs)=m(c-s)^2=mg>0.
 \label{eq:ch19-squeezed-check}
\end{equation}
The negative anomalous contraction is consistent with a positive
state: the \emph{complete} square remains nonnegative.

\section{Spatial correlations on a finite ring}
\label{sec:ch19-spatial-correlators}

The mode labels become positions by a finite Fourier sum.  Fix
$L=2h>0$, let $\Lambda=\mathbb Z/L\mathbb Z$, and set
\begin{equation}
 \zeta=e^{2\pi i/L},\qquad
 \chi_m(x)=\zeta^{mx},\qquad x\in\Lambda.
 \label{eq:ch19-spatial-character}
\end{equation}
These are characters of the cyclic lattice, so
$\chi_{-m}(x)=\overline{\chi_m(x)}$ and
$\chi_m(x)\chi_{-m}(y)=\chi_m(x-y)$.  The oscillator model remains
uncompressed in occupation number; only the positive mode set
$1\leq m\leq M$ is finite.  The algebraic calculation does not
require $M<h$, although a low-momentum identification with distinct
non-Nyquist lattice modes normally does.

Consider the \emph{nonzero-mode part} of the right-moving density,
\begin{equation}
 \rho_R(x)=\frac1L\sum_{m=1}^{M}
 \left(\chi_{-m}(x)C_{Rm}+\chi_m(x)A_{Rm}\right).
 \label{eq:ch19-spatial-rho}
\end{equation}
It is self-adjoint by the conjugation property of the characters.
This is not the full microscopic density: its charge zero mode is
not included.  The $1/L$ normalization is fixed by the Fourier
convention of the preceding chapters.

\begin{theorem}[Exact ordered spatial kernel]
\label{thm:ch19-spatial-ordered}
For $x,y\in\Lambda$, with $d=x-y$,
\begin{equation}
 \boxed{\displaystyle
 D_P(x,y):=\omega_P(\rho_R(x)\rho_R(y))
 =\frac1{L^2}\sum_{m=1}^{M}m
 \bigl(s^2\chi_{-m}(d)+c^2\chi_m(d)\bigr).}
 \label{eq:ch19-spatial-ordered}
\end{equation}
\end{theorem}
\begin{proof}
Expand the product of the two finite Fourier sums.  The $CC$ and
$AA$ expectations vanish.  In the $CA$ and $AC$ terms, the mode
contractions impose $m=n$ and provide, respectively, the factors
$ms^2$ and $mc^2$.  Multiplying the remaining characters gives
$\chi_{-m}(x-y)$ and $\chi_m(x-y)$, proving the formula.
\end{proof}

The kernel is generally \emph{complex}.  Let
$\theta_m=2\pi m(x-y)/L$, with any integer representatives of
$x,y$.  Using $c^2-s^2=1$ gives the particularly transparent form
\begin{equation}
 \boxed{\displaystyle
 D_P(x,y)=\frac1{L^2}\sum_{m=1}^{M}m
 \left[\frac{g+g^{-1}}2\cos\theta_m
                   +i\sin\theta_m\right].}
 \label{eq:ch19-spatial-ordered-realimag}
\end{equation}
The real part depends on the interacting state; the imaginary part
is fixed by the current algebra.  Taking the real part too early
would discard precisely the information distinguishing the two
operator orders.

\begin{proposition}[Commutator and symmetric fluctuations]
\label{prop:ch19-spatial-commutator}
\label{prop:ch19-spatial-symmetrized}
The abstract density operators obey
\begin{equation}
 [\rho_R(x),\rho_R(y)]
  =\frac{2i}{L^2}\sum_{m=1}^{M}m\sin\theta_m\,I,
 \label{eq:ch19-spatial-commutator}
\end{equation}
and their symmetrized two-point function is
\begin{equation}
 S_P(x,y):=\frac12\omega_P(\{\rho_R(x),\rho_R(y)\})
   =\frac{g+g^{-1}}{2L^2}
       \sum_{m=1}^{M}m\cos\theta_m.
 \label{eq:ch19-spatial-symmetrized}
\end{equation}
\end{proposition}
\begin{proof}
Compute the commutator directly from
$[A_{Rm},C_{Rn}]=m\delta_{mn}I$; only the mixed terms survive.
The symmetric formula is the half-sum of $D_P(x,y)$ and
$D_P(y,x)$.
\end{proof}

\subsection*{Positivity is a property of the whole kernel}

At coincident points the Fourier phases are one, so
\begin{theorem}[Coincident-point fluctuation]
\label{thm:ch19-spatial-diagonal}
\begin{equation}
 D_P(x,x)
 =\frac{c^2+s^2}{L^2}\frac{M(M+1)}2
 =\frac{g+g^{-1}}{2L^2}\frac{M(M+1)}2.
 \label{eq:ch19-spatial-diagonal}
\end{equation}
This is strictly positive for $M\geq1$ and vanishes for $M=0$.
\end{theorem}

But a positive two-point function need not be real or pointwise
nonnegative away from $x=y$.  The correct statement is matrix
positivity.
\begin{proposition}[Positive-semidefinite ordered kernel]
\label{prop:ch19-spatial-psd}
For sites $x_1,\ldots,x_r$ and coefficients $z_1,\ldots,z_r\in\mathbb C$,
\begin{equation}
 \sum_{a,b=1}^{r}\overline{z_a}D_P(x_a,x_b)z_b\geq0,
 \qquad D_P(y,x)=\overline{D_P(x,y)}.
 \label{eq:ch19-spatial-psd}
\end{equation}
\end{proposition}
\begin{proof}
With $B=\sum_b z_b\rho_R(x_b)$, the sum is exactly
$\omega_P(B^*B)\geq0$.  Hermiticity follows because each
$\rho_R(x)$ is self-adjoint.
\end{proof}

For a finite-mode theory this is already a complete result: the
Fourier kernel is explicit, cutoff dependent, Hermitian and positive
semidefinite.  We now show how vertex correlators and their characteristic
Luttinger exponents are constructed within this finite algebraic framework.

\subsection*{Finite-mode scalar Wightman correlators and the exponent dictionary}

In continuum bosonization, primary fields are expressed as normal-ordered exponentials
$:\!e^{i\alpha\phi}\!:$.  In our finite algebraic framework, we do not presuppose operator
convergence for infinite mode sums.  Instead, for any positive-mode cutoff $M \geq 1$ and
ultraviolet regulator length $a > 0$, we define normalized Hermitian algebraic field
linear combinations and deduce their scalar expectations.

Set $q_n = 2\pi n / L$ for $n \in \{1, \ldots, M\}$.  For chiralities $r \in \{+1, -1\}$
(with $r = +1$ for $R$ and $r = -1$ for $L$), define the finite-mode algebraic fields
\begin{equation}
 X_r^{(M)}(x) := \sum_{n=1}^M \frac{e^{-a q_n/2}}{\sqrt{n}}
 \left( e^{i r q_n x} b_{r,n} + e^{-i r q_n x} b_{r,n}^\dagger \right),
 \label{eq:ch19-normalized-field}
\end{equation}
where $b_{r,n} = C_{r,n}/\sqrt{n}$ and $b_{r,n}^\dagger = A_{r,n}/\sqrt{n}$ satisfy
standard CCR $[b_{r,n}, b_{s,m}^\dagger] = \delta_{rs}\delta_{nm} I$.
The bare vacuum two-point covariance is
\begin{equation}
 \omega_0\bigl(X_r^{(M)}(x) X_s^{(M)}(y)\bigr)
  = \delta_{rs} \sum_{n=1}^M \frac{e^{-a q_n}}{n} e^{i r q_n(x-y)},
 \label{eq:ch19-field-covariance}
\end{equation}
and their central commutator is
\begin{equation}
 [X_r^{(M)}(x), X_s^{(M)}(y)]
  = 2 i r \delta_{rs} \sum_{n=1}^M \frac{e^{-a q_n}}{n} \sin\bigl(q_n(x-y)\bigr) I.
 \label{eq:ch19-field-commutator}
\end{equation}

With the locked Bogoliubov coefficients $c, s$ from \eqref{eq:ch19-hyperbolic-locked},
the dressed chiral fields are
\begin{equation}
 A_R(x) := c X_R^{(M)}(x) + s X_L^{(M)}(x), \qquad
 A_L(x) := s X_R^{(M)}(x) + c X_L^{(M)}(x).
 \label{eq:ch19-dressed-chiral-fields}
\end{equation}
Because $c^2 - s^2 = 1$, the same-chirality commutator retains the canonical coefficient:
\begin{equation}
 [A_R(x), A_R(y)] = [X_R(x), X_R(y)].
\end{equation}
Moreover, the cross-chirality commutator vanishes identically at every finite $M$:
\begin{equation}
 [A_R(x), A_L(y)] = c s [X_R(x), X_R(y)] + s c [X_L(x), X_L(y)] = 0,
 \label{eq:ch19-cross-chiral-comm-zero}
\end{equation}
since the left branch acquires an extra minus sign from $r = -1$.

For real parameters $u, v \in \mathbb R$, consider the composite field
$Y_{u,v}(x) = u X_R^{(M)}(x) + v X_L^{(M)}(x)$.  Its vacuum expectation
value of the normal-ordered product defines the \emph{finite-mode scalar Wightman function}:

\begin{definition}[Finite-Mode Scalar Wightman Correlator]
\label{def:ch19-scalar-wightman}
For spatial separation $z = x - y$, define the scalar function
\begin{equation}
 \boxed{
 G_{u,v}^{(M)}(z)
  := \exp\left[ u^2 \sum_{n=1}^M \frac{e^{-a q_n}}{n} (e^{i q_n z} - 1)
               + v^2 \sum_{n=1}^M \frac{e^{-a q_n}}{n} (e^{-i q_n z} - 1) \right].
 }
 \label{eq:ch19-scalar-wightman}
\end{equation}
\end{definition}

\begin{proposition}[Modulus and Complex Phase of the Wightman Correlator]
\label{prop:ch19-wightman-modulus}
The modulus of $G_{u,v}^{(M)}(z)$ depends only on $u^2+v^2$ and the real cosine sum:
\begin{equation}
 \boxed{
 |G_{u,v}^{(M)}(z)|
  = \exp\left[ -(u^2+v^2) \sum_{n=1}^M \frac{e^{-a q_n}}{n} \bigl(1 - \cos(q_n z)\bigr) \right].
 }
 \label{eq:ch19-wightman-modulus}
\end{equation}
Its complex phase is $\exp[i (u^2 - v^2) \sum_{n=1}^M \frac{e^{-a q_n}}{n} \sin(q_n z)]$.
\end{proposition}
\begin{proof}
For any complex exponent $\exp(W)$, $|\exp(W)| = \exp(\operatorname{Re} W)$.
Here
$\operatorname{Re}(e^{i q_n z} - 1) = \cos(q_n z) - 1 = -(1 - \cos(q_n z))$,
and identically $\operatorname{Re}(e^{-i q_n z} - 1) = -(1 - \cos(q_n z))$.
Factoring $-(1 - \cos(q_n z))$ from both sums yields \eqref{eq:ch19-wightman-modulus}.
The imaginary part is $\operatorname{Im}(e^{\pm i q_n z} - 1) = \pm \sin(q_n z)$,
yielding the stated phase.
\end{proof}

\begin{theorem}[Continuum Limit and Conformal Power-Law Decay]
\label{thm:ch19-continuum-wightman}
At fixed ultraviolet cutoff $a > 0$, taking the infinite mode limit $M \to \infty$ yields
the convergent analytic expression
\begin{equation}
 G_{u,v}^{(\infty)}(z)
  = \left( \frac{1 - \rho}{1 - \rho e^{i t}} \right)^{u^2}
    \left( \frac{1 - \rho}{1 - \rho e^{-i t}} \right)^{v^2},
 \label{eq:ch19-wightman-continuum}
\end{equation}
where $\rho = e^{-2\pi a / L} \in (0, 1)$ and $t = 2\pi z / L$.
Consequently, the modulus is
\begin{equation}
 |G_{u,v}^{(\infty)}(z)|
  = \left[ \frac{1 - \rho}{\sqrt{1 + \rho^2 - 2\rho\cos t}} \right]^{u^2+v^2}.
 \label{eq:ch19-wightman-continuum-modulus}
\end{equation}
In the physical asymptotic regime $a \ll |z| \ll L$, this reproduces the power-law decay
\begin{equation}
 |G_{u,v}^{(\infty)}(z)| \approx \left( \frac{a}{|z|} \right)^{u^2+v^2}.
 \label{eq:ch19-asymptotic-power-law}
\end{equation}
\end{theorem}
\begin{proof}
Use the Taylor series $\sum_{n=1}^\infty \frac{w^n}{n} = -\ln(1 - w)$ for $|w| < 1$.
With $w = \rho e^{i t}$, the series converges absolutely because $\rho < 1$.
The first sum in \eqref{eq:ch19-scalar-wightman} becomes:
\[
 \sum_{n=1}^\infty \frac{\rho^n}{n} (e^{i n t} - 1)
  = \sum_{n=1}^\infty \frac{(\rho e^{i t})^n}{n} - \sum_{n=1}^\infty \frac{\rho^n}{n}
  = -\ln(1 - \rho e^{i t}) + \ln(1 - \rho)
  = \ln\left( \frac{1 - \rho}{1 - \rho e^{i t}} \right).
\]
Exponentiating with factor $u^2$, and repeating for the complex-conjugate second sum with factor $v^2$,
proves \eqref{eq:ch19-wightman-continuum}.
Taking the modulus:
$|1 - \rho e^{i t}|^2 = (1 - \rho\cos t)^2 + (\rho\sin t)^2 = 1 + \rho^2 - 2\rho\cos t$.
This proves \eqref{eq:ch19-wightman-continuum-modulus}.
For $a \ll L$ and $|z| \ll L$, $1 - \rho = 1 - e^{-2\pi a/L} \approx 2\pi a / L$,
and $1 + \rho^2 - 2\rho\cos t \approx (1-\rho)^2 + \rho t^2 \approx (2\pi/L)^2 (a^2 + z^2) \approx (2\pi z / L)^2$.
The quotient is $(a/|z|)$, raised to the power $u^2+v^2$.
\end{proof}

\subsection*{The Operator Exponent Dictionary}

Theorem~\ref{thm:ch19-continuum-wightman} establishes that the decay exponent $\gamma = u^2 + v^2$
is not a generic property of ``the vertex,'' but is uniquely determined by the weights $(u, v)$
of the chiral oscillator fields.  In terms of the dual bosonic fields
\begin{equation}
 \phi(x) = \frac{\sqrt{K}}{2}\bigl(X_R(x) + X_L(x)\bigr), \qquad
 \vartheta(x) = \frac{1}{2\sqrt{K}}\bigl(X_R(x) - X_L(x)\bigr),
 \label{eq:ch19-dual-fields-def}
\end{equation}
we obtain the exact exponent dictionary:

\begin{theorem}[Operator Exponent Dictionary]
\label{thm:ch19-exponent-dictionary}
Under Bogoliubov dressing with Luttinger parameter $K > 0$, the physical order parameters
and elementary fermion fields possess the following equal-time decay exponents:
\begin{enumerate}
 \item \textbf{Charge-Density-Wave (CDW) Harmonic:}\\
       The density oscillation operator $\mathcal O_{\mathrm{CDW}}(x) \sim :\!e^{2i\phi(x)}\!:$ has
       chiral weights $u = v = \sqrt{K}$.  Its decay exponent is
       \begin{equation}
        \gamma_{\mathrm{CDW}} = u^2 + v^2 = K + K = 2K.
        \label{eq:ch19-cdw-exponent}
       \end{equation}
 \item \textbf{Superconducting (SC) Pairing:}
       The pairing operator $\mathcal O_{\mathrm{SC}}(x) \sim :\!e^{2i\vartheta(x)}\!:$ has
       chiral weights $u = 1/\sqrt{K}, v = -1/\sqrt{K}$.  Its decay exponent is
       \begin{equation}
        \gamma_{\mathrm{SC}} = u^2 + v^2 = \frac{1}{K} + \frac{1}{K} = \frac{2}{K}.
        \label{eq:ch19-sc-exponent}
       \end{equation}
 \item \textbf{Elementary Dressed Fermion:}
       The right-moving dressed fermionic vertex $\Psi_R(x) \sim :\!e^{i A_R(x)}\!:$ has
       chiral weights $u = c, v = s$.  Its decay exponent is
       \begin{equation}
        \gamma_{\Psi} = c^2 + s^2 = \frac{K + K^{-1}}{2}.
        \label{eq:ch19-fermion-exponent}
       \end{equation}
\end{enumerate}
\end{theorem}

\begin{table}[ht]
\centering
\caption{The Operator Exponent Dictionary across physical regimes.}
\label{tab:ch19-exponent-dictionary}
\renewcommand{\arraystretch}{1.3}
\resizebox{\textwidth}{!}{%
\begin{tabular}{lcccc}
\hline\hline
\textbf{Operator} & \textbf{Bosonic Representation} & \textbf{General $K$} & \textbf{Free ($K=1$)} & \textbf{Repulsive ($K < 1$)} \\
\hline
CDW density harmonic & $:\!e^{2i\phi(x)}\!:$ & $2K$ & $2$ & $< 2$ (enhanced) \\
SC pairing order & $:\!e^{2i\vartheta(x)}\!:$ & $2/K$ & $2$ & $> 2$ (suppressed) \\
Elementary chiral fermion & $:\!e^{i A_R(x)}\!:$ & $\frac{K+K^{-1}}{2}$ & $1$ & $> 1$ (suppressed) \\
\hline\hline
\end{tabular}%
}
\end{table}

At the non-interacting point $K = 1$, the CDW and superconducting exponents both equal $2$,
while the free fermion Green function has exponent $1$, reproducing the standard Fermi liquid
and free Dirac behavior in $1+1$ dimensions.  For repulsive interactions ($K < 1$), CDW fluctuations
decay more slowly ($\gamma_{\mathrm{CDW}} < 2$) and dominate over superconductivity ($\gamma_{\mathrm{SC}} > 2$),
while the single-particle fermion propagator decays faster than $1/|z|$, signalling the destruction
of quasiparticle poles and the emergence of a Tomonaga--Luttinger liquid.

\section{Returning to the finite fermionic sea}
\label{sec:ch19-finite-comparison}

We now evaluate \emph{corresponding formulas} in the original finite
fermionic representation.  The word ``corresponding'' matters: the
abstract generators and the finite CAR currents act on different
spaces.  Equality of two calculated numbers need not provide a map
identifying their operator algebras or states.

\subsection*{Protected current contractions}

Fix $L=2h$, a charge assignment $N$ admissible for the finite
construction, and its normalized bare sector-ground vector $\Omega_N$.
Write
\begin{equation}
 R_m:=\rho^{\mathrm F}_{m,R},\qquad
 R_{-m}:=\rho^{\mathrm F}_{-m,R}=R_m^*,
 \qquad \langle X\rangle^{\mathrm F}_N
 :=\langle\Omega_N,X\Omega_N\rangle.
 \label{eq:ch19-comp-two-evaluators}
\end{equation}
Assume the right-branch \emph{zero-budget margin} used in the finite
current theorem, together with $h>0$ and $1\leq m,n\leq M$.
In the source formalization this is
$\mathrm{marginAt}(h,M,0,N,R)$; its excursion bounds include
$m+n+|N_R|\leq h$.  The relevant vector identities are
\begin{equation}
 R_{-m}\Omega_N=0,
 \qquad [R_{-m},R_n]\Omega_N=m\delta_{mn}\Omega_N.
 \label{eq:ch19-comp-ground-commutator}
\end{equation}
They are only asserted on the specified state under these
hypotheses, not on the entire finite CAR Fock space.

\begin{proposition}[Protected finite-ground contractions]
\label{prop:ch19-comp-ground-contractions}
For the modes in this range,
\begin{align}
 \langle R_{-m}R_n\rangle^{\mathrm F}_N&=m\delta_{mn},
   \label{eq:ch19-comp-lr}\\
 \langle R_mR_{-n}\rangle^{\mathrm F}_N&=0,
   \label{eq:ch19-comp-rl}\\
 \langle R_mR_n\rangle^{\mathrm F}_N
  &=\langle R_{-m}R_{-n}\rangle^{\mathrm F}_N=0.
   \label{eq:ch19-comp-same-sign}
\end{align}
In particular, $\|R_m\Omega_N\|^2=m$.
\end{proposition}
\begin{proof}
Use
$R_{-m}R_n=[R_{-m},R_n]+R_nR_{-m}$ on $\Omega_N$.
The second term vanishes because the lowering mode acts first.
The reversed product and the two-lowering product vanish for the
same reason.  The two-raising expectation vanishes by adjointness,
and the norm formula follows from the first contraction.
\end{proof}

\subsection*{The precise free-state agreement}

Build the right-branch finite oscillatory density with the same
selected modes and Fourier convention:
\begin{equation}
 \rho_R^{\mathrm F}(x)
 :=\frac1L\sum_{m=1}^{M}
 \left(\chi_{-m}(x)R_m+\chi_m(x)R_{-m}\right).
 \label{eq:ch19-comp-fock-density}
\end{equation}

\begin{theorem}[Free quadratic comparison]
\label{thm:ch19-comp-free-kernel}
Under the finite current-margin hypotheses above, for all
$x,y\in\Lambda$,
\begin{equation}
 \boxed{\displaystyle
 \langle\rho_R^{\mathrm F}(x)\rho_R^{\mathrm F}(y)
 \rangle^{\mathrm F}_N
 =\frac1{L^2}\sum_{m=1}^{M}m\chi_m(x-y)
 =\omega_0(\rho_R(x)\rho_R(y)).}
 \label{eq:ch19-comp-fock-kernel}
\end{equation}
\end{theorem}
\begin{proof}
Expand the ordered product of the two finite densities.
The protected contractions kill every term except
$R_{-m}R_n$, which contributes $m\delta_{mn}$.  The remaining
characters give the displayed sum.  For the abstract state,
substitute the free hyperbolic pair $(c,s)=(1,0)$ into
\eqref{eq:ch19-spatial-ordered}.
\end{proof}

The agreement is exact for these quadratic observables.  It does
not identify the finite and abstract states on arbitrary words,
nor does it identify the bare sea with an interacting ground state.

\subsection*{What changes when the mixing is nonzero?}

For the two chosen finite branches let
\begin{equation}
 \widetilde A_{\nu m}^{\mathrm F}
   =c\rho_{-m,\nu}^{\mathrm F}
     +s\rho_{m,\bar\nu}^{\mathrm F},
 \qquad
 \widetilde C_{\nu m}^{\mathrm F}
   =(\widetilde A_{\nu m}^{\mathrm F})^*.
 \label{eq:ch19-comp-dressed-ann}
\end{equation}
Assume the \emph{both-branch} zero-budget margins, so that the
relevant excitation norms hold for both $R$ and $L$.  Since the
bare lowering currents kill $\Omega_N$,
$\widetilde A_{\nu m}^{\mathrm F}\Omega_N
=s\rho_{m,\bar\nu}^{\mathrm F}\Omega_N$.

\begin{proposition}[Finite dressed-number witness]
\label{prop:ch19-comp-dressed-number}
Under these assumptions, for $1\leq m\leq M$,
\begin{equation}
 \langle\widetilde C_{\nu m}^{\mathrm F}
           \widetilde A_{\nu m}^{\mathrm F}\rangle_N^{\mathrm F}
   =ms^2,
 \qquad
 \langle\widetilde A_{\nu m}^{\mathrm F}
           \widetilde C_{\nu m}^{\mathrm F}\rangle_N^{\mathrm F}
   =mc^2.
 \label{eq:ch19-comp-dressed-ca}
\end{equation}
If $s\ne0$, the bare sea $\Omega_N$ is not a vacuum for the
finite dressed annihilator.
\end{proposition}
\begin{proof}
The first expectation is
$\|\widetilde A_{\nu m}^{\mathrm F}\Omega_N\|^2
=s^2\|\rho_{m,\bar\nu}^{\mathrm F}\Omega_N\|^2=ms^2$.
For the second, use
$\widetilde C_{\nu m}^{\mathrm F}\Omega_N
=c\rho_{m,\nu}^{\mathrm F}\Omega_N$ and take its squared norm.
\end{proof}

The abstract \emph{bare} contraction
$\omega_P(C_{\nu m}A_{\nu m})$ and the finite \emph{dressed}
contraction
$\langle\widetilde C_{\nu m}^{\mathrm F}\widetilde A_{\nu m}^{\mathrm F}\rangle_N^{\mathrm F}$
both equal $ms^2$, but refer to different operators in different
states.  In fact,
$\omega_P(\widetilde C_{\nu m}\widetilde A_{\nu m})=0$,
while the finite value is nonzero for $s\ne0$.

\begin{example}[Sixteen-site finite/abstract comparison]
\label{ex:ch19-comp-16-sites}

Use the same stable hyperbolic pair that appears in the Chapter~18
finite witness:
\[
 L=16,\qquad h=8,\qquad M=2,\qquad
 c=\frac54,\quad s=\frac34,\quad g=(c-s)^2=\frac14.
\]
Take the admissible zero-charge sea, for which the stated
both-branch margins hold.  At the second positive mode,
\begin{equation}
 \|R_2\Omega_0\|^2=2,
 \qquad
 \langle\widetilde C_{R2}^{\mathrm F}
          \widetilde A_{R2}^{\mathrm F}\rangle_0^{\mathrm F}
 =2\left(\frac34\right)^2=\frac98.
 \label{eq:ch19-comp-example-dressed}
\end{equation}
At one site the \emph{finite bare} density variance is
\begin{equation}
 \langle(\rho_R^{\mathrm F}(0))^2\rangle_0^{\mathrm F}
   =\frac{1+2}{16^2}=\frac3{256},
 \label{eq:ch19-comp-example-free}
\end{equation}
while the \emph{abstract interacting} density variance is
\begin{equation}
 D_P(0,0)=\frac{1+2}{16^2}
 \left(\frac{25}{16}+\frac9{16}\right)
 =\frac{51}{2048}\ne\frac3{256}.
 \label{eq:ch19-comp-example-abstract}
\end{equation}
The comparison is deliberately between a \emph{bare} finite sea
and an \emph{interacting} abstract state.  It neither constructs nor
rules out a different finite interacting ground state.
\end{example}

For a separate illustration of the danger of energy truncation,
retain only the three polynomial states
$\{1,X_R,X_L\}$ of total degree at most one.  Compressing the
two dressed annihilators gives
\[
 Q_R=\begin{pmatrix}0&c&0\\0&0&0\\s&0&0\end{pmatrix},
 \qquad
 Q_L=\begin{pmatrix}0&0&c\\s&0&0\\0&0&0\end{pmatrix},
 \qquad
 Q_R^*Q_R+Q_L^*Q_L
   =\operatorname{diag}(2s^2,c^2,c^2).
\]
For $s\ne0$ the diagonal matrix is strictly positive, so
$\ker Q_R\cap\ker Q_L=\{0\}$.  Nevertheless the positive
finite Hamiltonian $Q_R^*Q_R+Q_L^*Q_L$ has ground states.
Annihilation-vacuum conditions and ground-state existence are
not equivalent assertions.

\subsection*{The remaining comparison problem}

To compare the interacting abstract state $\omega_P$ with a genuine
interacting finite-CAR state, one must first choose the latter
(e.g.\ a ground density matrix $\Gamma_h$ for a specified finite
Hamiltonian), and choose which abstract words $a$ are represented
by which finite operators $\Phi_h(a)$.  For a given word, the
comparison defect can then be stated precisely as
\begin{equation}
 \varepsilon_h(a)
  :=\operatorname{Tr}(\Gamma_h\Phi_h(a))-\omega_P(a).
 \label{eq:ch19-comp-comparison-defect}
\end{equation}
The map $\Phi_h$ cannot be a globally exact representation of the
nonzero scalar CCR by finite matrices.  A meaningful comparison must
specify its domain of words and control intermediate excursions,
finite-band edges, and the chosen energy budgets.

\begin{theorem}[Word-Dependent Protected Bridge]
\label{thm:ch19-protected-bridge}
Let $\mathcal W$ be a finite family of current words in the bare density generators.
Let $\mathcal D_B$ be a finite-dimensional subspace of bosonic excitation polynomials,
and let $\iota\colon \mathcal D_B \hookrightarrow \mathcal H_{\mathrm F}$ be an isometric
embedding into the finite fermionic Fock space.
Suppose the certified band-margin condition $\mathrm{marginAt}(h, M, \mathrm{budget})$
holds for all intermediate states reached during the action of every word $w \in \mathcal W$.
Then, for all $u, v \in \mathcal D_B$ and every $w \in \mathcal W$:
\begin{equation}
 \boxed{\langle\iota u, \pi_{\mathrm F}(w) \iota v\rangle_{\mathrm F}
  = \langle u, \pi_{\mathrm B}(w) v\rangle_{\mathrm B}.}
 \label{eq:ch19-word-bridge}
\end{equation}
\end{theorem}
\begin{proof}
By the margin hypothesis, the action of every generator in the word $w$ leaves the
fermionic state inside the protected domain where the boundary corrections
$\sum_{k=-B}^{-B+n-1} \hat{n}_k - \sum_{k=B-n+1}^B \hat{n}_k$ act strictly as the central
scalar $n\cdot I$ (Theorem~\ref{thm:ch19-current-boundary-commutator} and
Corollary~\ref{cor:ch19-margin-protected-action}).
Because no intermediate projection is inserted (avoiding the compression remainder
\eqref{eq:ch19-compression-remainder}), the uncompressed operator product intertwines with
$\iota$ on $\mathcal D_B$, reducing the matrix element to the evaluation in the polynomial
CCR representation.
\end{proof}

The significance of Theorem~\ref{thm:ch19-protected-bridge} is architectural: it proves
matrix-element equality for \emph{certified words on protected domains}, rather than
postulating a false global operator isomorphism $\mathcal B(\mathcal H_{\mathrm F}) \cong \mathfrak A_M$.

The outcome of this chapter is therefore sharp but limited.  We have
constructed a positive algebraic dressed state; calculated its ordered
mode and spatial correlations; established the positivity of the
resulting kernel; derived the scalar Wightman function and the Operator Exponent Dictionary;
and matched its free quadratic kernel and protected words to a finite fermionic sea.
The Gaussian state is not automatically the finite interacting ground state.
Chapter~20 uses this distinction when it introduces genuine finite-CAR products and, separately,
formal vertex correlators.  Neither a convergent vertex exponential
nor a continuum power law has been derived from the present
finite-mode results without the explicit hypotheses stated above.

\paragraph{Verification boundary.}
CH19/A08 Phase C is complete within its reviewed scope: the polynomial
oscillator state, ordered mode contractions, spatial kernel, scalar Wightman
correlators, and protected free finite-state comparison are promoted and audited in Core.  The
independent validation record is\footnote{\texttt{note/ch19\_phase\_c\_validation\_2026-10-16.md}}.
This chapter remains a manuscript candidate pending its own independent
editorial review.  The Lean result does not establish a finite interacting
ground-state equivalence, analytic evaluation of formal vertices, or a
continuum power law; those remain separate obligations.
